\documentclass[12pt]{report}
\usepackage[utf8]{inputenc}
\usepackage[spanish]{babel}
\usepackage{textcomp}
\usepackage{graphicx}
\usepackage{amsmath}
\usepackage{geometry}
\usepackage{colortbl}
\usepackage{array}
\usepackage{float}
\usepackage{enumitem}
\geometry{left=2.5cm, right=2.5cm, top=2.5cm, bottom=2.5cm}
\usepackage{titlesec}

\bibliographystyle{unsrt}


% Estilo para títulos de sección
\titleformat{\section}{\normalfont\Large\bfseries}{\thesection}{1em}{}
\usepackage{atbegshi}
\spanishdecimal{.}
\begin{document}

% Portada estilizada

    \begin{center}
    
        % Logos e institución
        
\includegraphics[width=4cm]{Reporte Semestre 1/logos/uaem.png} \hfill
        
\includegraphics[width=5cm]{Reporte Semestre 1/logos/ciicap.png} \\[2cm]
        
        % Línea separadora superior
        \hrulefill \\[0.4cm]

        % Programa educativo
        {\Large \textsc{Maestría en Ingeniería y Ciencias Aplicadas}} \\[0.5cm]

        % Título del reporte
        \hrulefill \\[0.4cm]
        {\LARGE \textbf{Fluorescencia Resonante Intermitente con Ecuaciones Estocásticas de Schrodinger}} \\[0.4cm]
        \hrulefill \\[0.5cm]

        % Subtítulo y periodo
        {\Large \textbf{Tercer Reporte Semestral}} \\[0.2cm]
        {\large \textbf{Periodo: Agosto -Diciembre 2025}} \\[2cm]
        
        \vfill

        % Información del estudiante
        \textbf{Alumno:} \\[0.2cm]
        {\Large L. T. Kevin Martínez Franco} \\[1.5cm]

        % Información del asesor
        \textbf{Asesor:} \\[0.2cm]
        {\Large Dr. Héctor Manuel Castro Beltrán} \\[2cm]

        % Fecha
        {\large \textbf{\today}} \\[2cm]
        
        % Línea final
        \hrulefill

    \end{center}


% Resumen
\section*{Resumen}

Durante este semestre se avanzó significativamente en el estudio de la fluorescencia resonante intermitente, obteniendo resultados nuevos en varios frentes. En primer lugar, se optimizaron los espectros de detección heterodina balanceada, logrando una reducción notable del ruido y permitiendo una comparación directa y consistente con los espectros obtenidos mediante la ecuación maestra. En segundo lugar, se implementó la detección homodina balanceada; las cuadraturas en $\pi/2$ y 0 permitieron identificar, respectivamente, el pico angosto con laterales y el pico central coherente, cuya combinación reconstruye el espectro completo. En tercer lugar, se desarrolló un análisis fotoestadístico de los tiempos de retorno al estado metastable, aplicando teoría de procesos de renovación. Al inicializar en el estado metastable y registrar los tiempos de reocupación, se obtuvieron distribuciones que evidencian la alternancia entre periodos brillantes y oscuros, con predicciones teóricas que se confirmaron mediante simulaciones tanto en promedios como en momentos. Finalmente, se generó una pseudo-esfera de Bloch para el sistema de tres niveles, permitiendo visualizar la dinámica estocástica y contrastar directamente la evolución en los métodos de saltos cuánticos y trayectorias difusivas.

En conjunto, estos avances fortalecen el enfoque general de la tesis, orientado a caracterizar de manera sistemática la fluorescencia intermitente mediante ecuaciones estocásticas de Schrödinger, integrando detección homodina, heterodina y análisis fotoestadístico para describir de forma clara el comportamiento del sistema.



% Introducción
\section*{Introducción}


La disipación, la pérdida irreversible de energía y coherencia de un microsistema, es el resultado del acoplamiento a un macrosistema (o reservorio) mucho más grande, que es tan grande que no hay manera de llevar un registro de todos sus grados de libertad.

En un sistema cuántico abierto, la descripción estándar del ensamble es la ecuación maestra, que gobierna la evolución del operador densidad y captura la disipación y el ruido inducidos por el entorno. Sin embargo, la ecuación maestra describe solo el comportamiento promedio de muchos sistemas idénticos; no especifica cómo evoluciona un solo sistema cuántico condicionado a una medición continua.

Para conectar esa dinámica promedio con una dinámica ``por trayectoria'', se introducen las ecuaciones estocásticas de Schrödinger (SSE) o ecuaciones de evolución condicionada. Existen infinitas descomposiciones estocásticas de una misma ecuación maestra (los llamados \textit{unravelings}): cada una consiste en una dinámica aleatoria para un estado puro o mixto tal que, al promediar sobre todas las realizaciones estocásticas, se recupera exactamente la ecuación maestra \cite{Carmichael1993,WisemanMilburn2010}.

Pero ---y aquí está el punto central--- aunque matemáticamente hay una infinidad de \textit{unravelings}, solo algunos corresponden a esquemas de medición físicamente realizables. Es decir, solo ciertas SSE describen cómo evoluciona el estado del sistema condicionado a un tipo específico de detección experimental, como detección de conteo de fotones (saltos cuánticos), detección homodina o heterodina (trayectorias difusivas), etc. \cite{Carmichael1993}



Históricamente, las trayectorias estocásticas ofrecían una ventaja computacional clara frente a la integración de la ecuación maestra, debido a que la evolución de un estado puro de dimensión $d$ es mucho menos costosa que la evolución de un operador densidad de dimensión $d^2$.  
Aunque el aumento en la capacidad de cómputo moderna ha reducido esta diferencia para sistemas pequeños, existen actualmente numerosos casos donde las trayectorias siguen siendo superiores: sistemas de muchos cuerpos o de gran dimensión de Hilbert \cite{Daley2014,Lindner2022}, simulaciones paralelizables en GPU \cite{Li2022}, sistemas bosónicos o de circuit QED con grandes truncamientos \cite{Koch2019}, y dinámicas no-Markovianas donde la ecuación maestra resulta impracticable \cite{Suess2014,Ke2021}.

No obstante, para los propósitos del presente proyecto, la motivación principal para emplear trayectorias no es su posible eficiencia computacional, sino su \textbf{interpretación física como estados condicionados por resultados de medición}. Esta perspectiva proporciona una intuición directa sobre la relación entre la dinámica irreversible del sistema y los registros continuos o discretos del entorno, lo cual resulta central para el análisis conceptual desarrollado en este trabajo.



\section*{Trayectorias de Saltos Cuánticos y Distribuciones de Tiempos de Espera}
\label{sec:quantum_jumps}

El formalismo de trayectorias de saltos cuánticos proporciona una descripción estocástica de la evolución temporal de sistemas cuánticos abiertos, donde la disipación se manifiesta como transiciones discretas entre estados \cite{Carmichael1993,Plenio1998}. Este enfoque es particularmente relevante para sistemas atómicos individuales donde es posible observar directamente los saltos cuánticos, como en el efecto de la Flourosencia Resonante Intermitente.

\subsection*{Descomposición de la Ecuación Maestra}

La ecuación maestra para un sistema cuántico abierto puede escribirse como:
\begin{equation}
\dot{\rho} = \mathcal{L}\rho = -i(H\rho - \rho H^\dagger) + \sum_k \mathcal{D}[C_k]\rho ,
\label{eq:master_eq}
\end{equation}
donde $\mathcal{L}$ es el superoperador de Lindblad, 
\begin{equation}
H = H_0 - \frac{i}{2} \sum_k C_k^\dagger C_k
\label{eq:effective_hamiltonian}
\end{equation}
es un Hamiltoniano efectivo no Hermitiano, con $H_0$ siendo el Hamiltoniano del sistema en la aproximación de onda rotante (RWA), que para un átomo de tres niveles bajo driving láser toma la forma:
\begin{equation}
H_0 = -\hbar\Delta A_{ee} + \frac{\hbar\Omega}{2}(A_{eg} + A_{ge}),
\end{equation}
donde $\Delta = \omega_L - \omega_0$ es la desintonía láser-átomo y $\Omega$ es la frecuencia de Rabi. El término no-hermitiano $-\frac{i}{2} \sum_k C_k^\dagger C_k$ captura la disipación por todos los canales de decaimiento. Los operadores de disipación
\begin{equation}
\mathcal{D}[C_k]\rho = C_k\rho C_k^\dagger - \frac{1}{2}(C_k^\dagger C_k\rho + \rho C_k^\dagger C_k)
\label{eq:dissipation_operator}
\end{equation}
corresponden a eventos de emisión espontánea que colapsan el sistema a los estados correspondientes.

Para el sistema de tres niveles, los operadores de colapso son:
\begin{align}
C_{eg} &= \sqrt{\gamma} A_{ge}, \\
C_{ea} &= \sqrt{\gamma_d} A_{ae}, \\
C_{ag} &= \sqrt{\gamma_a} A_{ga},
\end{align}
que representan los tres canales de decaimiento posibles.

\subsection*{Interpretación de Trayectorias}

La solución formal de la ecuación maestra $\rho(t) = e^{\mathcal{L}t}\rho(0)$ puede descomponerse como:
\begin{equation}
\rho(t) = \sum_{m=0}^{\infty} \int_0^t dt_m \int_0^{t_m} dt_{m-1} \cdots \int_0^{t_2} dt_1 e^{(\mathcal{L}-\mathcal{D})(t-t_m)}
\mathcal{D}e^{(\mathcal{L}-\mathcal{D})(t_m-t_{m-1})}\mathcal{D} \cdots \mathcal{D}e^{(\mathcal{L}-\mathcal{D})t_1}\rho(0) ,
\label{eq:path_integral}
\end{equation}
donde $\mathcal{D} = \sum_k \mathcal{D}[C_k]$ es el superoperador de disipación total.

Cada término en esta serie representa una historia única de emisiones desde $t=0$ hasta $t$, donde $m$ es el número total de emisiones y $\{t_1, t_2, \ldots, t_m\}$ son los tiempos de emisión ordenados. El operador de densidad condicional no normalizado está dado por:
\begin{equation}
\bar{\rho}_c(t) = e^{(\mathcal{L}-\mathcal{D})(t-t_m)}\mathcal{D}e^{(\mathcal{L}-\mathcal{D})(t_m-t_{m-1})}\mathcal{D} \cdots \mathcal{D}e^{(\mathcal{L}-\mathcal{D})t_1}\rho(0) ,
\label{eq:conditional_density}
\end{equation}
y el operador de densidad condicional normalizado es:
\begin{equation}
\rho_c(t) = \frac{\bar{\rho}_c(t)}{\text{Tr}[\bar{\rho}_c(t)]}.
\label{eq:normalized_density}
\end{equation}

\subsection*{Distribuciones de Tiempos de Espera y el Fenómeno de la Fluorescencia Resonante Intermitente}

Las distribuciones de tiempos de espera caracterizan la estadística temporal entre eventos de emisión consecutivos, revelando correlaciones cuánticas como el antibunching. En sistemas que exhiben el fenómeno de \textit{Fluorescencia Resonante Intermitente} o shelving, estas distribuciones se manifiestan como secuencias alternantes de \textit{tiempos brillantes} (períodos de emisión intensa) y \textit{tiempos oscuros} (períodos de nula emisión) \cite{Knight1990}.

Para un átomo de dos niveles, la distribución de tiempos de espera viene dada por:
\begin{equation}
w(\tau) = \gamma \langle \sigma_+\sigma_-\rangle_{ss} e^{-\gamma \langle \sigma_+\sigma_-\rangle_{ss} \tau},
\label{eq:waiting_time_2level}
\end{equation}
pero cuando se considera la dinámica condicionada a no-detecciones, se obtiene la distribución correcta:
\begin{equation}
w(\tau) = \gamma |\langle \sigma_-\rangle(\tau)|^2 e^{-\gamma \int_0^\tau |\langle \sigma_-\rangle(\tau')|^2 d\tau'},
\label{eq:waiting_time_conditional}
\end{equation}
donde $\langle \sigma_-\rangle(\tau)$ es la evolución del operador de dipolo condicionada a no haber detecciones.

En la Florescencia Intermitente, la dinámica de atrapamiento en estados metaestables genera estadísticas de parpadeo observables experimentalmente, donde los tiempos brillantes corresponden a la emisión desde el estado excitado y los tiempos oscuros al atrapamiento en el estado metaestable \cite{Wiseman1993}.

\subsection*{Interpretación de Resultados}

La forma de $w(\tau)$ revela información crucial sobre la estadística del proceso de emisión:

\begin{itemize}
\item \textbf{Emisión Poissoniana:} Para una fuente coherente ideal, $w(\tau) = \gamma e^{-\gamma\tau}$, mostrando una decadencia exponencial simple.

\item \textbf{Antibunching:} $w(0) = 0$ indica que los fotones tienden a emitirse separados en tiempo, característico de sistemas como átomos individuales o puntos cuánticos.

\item \textbf{Bunching:} $w(0) > w(\tau)$ para $\tau$ pequeños sugiere que los fotones tienden a emitirse en grupos.
\end{itemize} 
\begin{figure}[h]
    \centering
    \includegraphics[width=0.8\linewidth]{Reporte Semestre 3/imagenes/Trayectoria de saltos.png}
    \caption{(a) Trayectoria de saltos cuánticos mostrando periodos brillantes y oscuros en un sistema de 3 niveles con $\Omega = 3.5\gamma$ y $\Delta t=0.001$. }
    \label{fig:quantum_jumps_shelving}
\end{figure}



\subsection*{Aplicación a la Floursecencia Intermitente}

El formalismo de los tiempos de espera puede extenderse para caracterizar fenómenos más complejos, como el parpadeo cuántico o intermitencia, donde un emisor alterna aleatoriamente entre periodos de emisión brillante 
y periodos de emisión oscura. El problema del \"shelving\" o almacenamiento en un estado metaestable es un mecanismo físico que da lugar a este comportamiento.




El procedimiento de c\'alculo de trayectorias permite recoger la informaci\'on de fotoemisiones y estad\'istica de los per\'iodos brillantes y oscuros[3]. Esto se logra mejor si hacemos una trayectoria muy larga en lugar de muchas trayectorias cortas. Se pueden obtener las distribuciones de las duraciones de los periodos oscuros y brillantes, que siguen distribuciones exponenciales: 
$$ P_O = \frac{1}{T_O} e^{-t/T_O}, \qquad  
    P_B = \frac{1}{T_B} e^{-t/T_B} ,$$
con duraciones promedio
$$ T_O = \gamma_a^{-1}, \qquad T_B = \frac{2\Omega^2 + \gamma^2 + 4\Delta^2}{\gamma_d \Omega^2} .$$
\\

\begin{itemize}
    \item \textbf{Distribución de Duración de Periodos Brillantes ($P_O(t)$):} Mide la estadística de la duración de los periodos de alta tasa de emisión. Corresponde a los tiempos entre una serie de saltos del estado excitado al estado base y termina cuando entra al estado metaestable.

    \item \textbf{Distribución de Duración de Periodos Oscuros (\(P_O(t)\)):} Mide la estadística de la duración de los periodos de nula emisión de fotones. Corresponde a los tiempos de permanencia en el estado metaestable, que finalizan con un salto de regreso al estado base.
    
    \item \textbf{Distribución de Regreso al Estado Metaestable (\(P_{OB}(t)\)):} Esta distribución caracteriza la duración total de un ciclo completo de parpadeo. Se obtiene al medir la suma de un tiempo oscuro (\(t_O\)) y el tiempo brillante (\(t_B\)) inmediatamente siguiente, usando como condición inicial el estado metaestable. No es una distribución de tiempos de espera en el sentido estricto, sino la distribución de la \textbf{suma de dos variables aleatorias independientes}. Su forma analítica, dada por la convolución de \(P_O(t)\) y \(P_B(t)\), es:
    $$
     P_{OB}(t) =   \dfrac{1}{T_B - T_O} \left( e^{-t/T_B} - e^{-t/T_O} \right)
    $$
  

    El tiempo promedio para completar este ciclo es simplemente la suma de los promedios individuales:
    \[
    T_{OB} = T_O + T_B.
    \]
\end{itemize}

\section*{Detección Homodina Balanceada}

La detección homodina balanceada (DHB) es una técnica fundamental para la medición de ruido cuántico y, en particular, para la caracterización de \textit{estados comprimidos de la luz}. A diferencia de la detección heterodina, en este esquema el oscilador local (OL) tiene la misma frecuencia que el campo bajo estudio ($\omega_{\text{lo}} = \omega$). Ambos haces, el del emisor y el del OL, suelen provenir del mismo láser. La clave de este método reside en la capacidad de seleccionar la fase relativa $\phi$ del OL, lo que permite medir la varianza de una cuadratura específica del campo.

\subsection*{Principio de Operación}

El esquema de la DHB (ver Fig. \ref{fig:homodina_balanceada}) consiste en mezclar el modo del campo óptico de interés, $\hat{a}$, con el modo de un OL intenso y coherente, $\hat{b}$, en un divisor de haz (DH) balanceado ($R = T = 1/2$). Los modos de salida tras el divisor de haz están dados por:

\begin{eqnarray}
\hat{c} = \frac{1}{\sqrt{2}} (\hat{a} + i\hat{b}), \quad
\hat{d} = \frac{1}{\sqrt{2}} (i\hat{a} + \hat{b})
\end{eqnarray}

Las intensidades registradas en los fotodetectores D1 y D2 son proporcionales a $\hat{c}^\dagger \hat{c}$ y $\hat{d}^\dagger \hat{d}$, respectivamente. La esencia de la técnica ``balanceada'' es restar estas dos fotocorrientes para obtener la señal diferencia:

\begin{eqnarray}
\hat{n}_{cd} = \hat{c}^\dagger \hat{c} - \hat{d}^\dagger \hat{d} = -i(\hat{a}^\dagger \hat{b} - \hat{b}^\dagger \hat{a})
\label{eq:diferencia_corrientes}
\end{eqnarray}

Esta operación cancela el exceso de ruido clásico del OL, cuya intensidad es típicamente mucho mayor que la del emisor, permitiendo acceder a las fluctuaciones cuánticas del campo $\hat{a}$.

\subsection{Sensibilidad a las Cuadraturas}

Considerando que el OL se encuentra en un estado coherente de gran amplitud, $\hat{b} = \beta e^{i\phi}$, podemos aproximar $\hat{b}$ como un número complejo $\beta e^{i\phi}$. Sustituyendo en la Ecuación \ref{eq:diferencia_corrientes}, la señal medida se simplifica a:

\begin{eqnarray}
\hat{n}_{cd} = i\beta\left(\hat{a}^\dagger e^{i\phi} - \hat{a} e^{-i\phi}\right)
\end{eqnarray}

Manipulando algebraicamente esta expresión, se puede demostrar que es proporcional a un operador de cuadratura:

\begin{eqnarray}
\hat{n}_{cd} = 2\beta \hat{X}_{\phi + \pi/2} = 2\beta \hat{P}_{\phi}
\end{eqnarray}

donde $\hat{X}_{\theta} = \frac{1}{2}(\hat{a}^\dagger e^{i\theta} + \hat{a} e^{-i\theta})$ es el operador de cuadratura generalizado. Al normalizar por la amplitud del OL, la señal $\hat{n}_{cd}$ proporciona una medición directa de la cuadratura $\hat{P}_{\phi}$, que es la conjugada de $\hat{X}_{\phi}$. Variando la fase $\phi$ del OL, es posible medir cualquier cuadratura del campo, permitiendo así caracterizar completamente sus propiedades de ruido.

\subsection{Aplicaciones: Detección de Compresión}

La DHB es la técnica por excelencia para detectar estados comprimidos de la luz. Un estado está comprimido si la varianza de una de sus cuadraturas, $\Delta^2 \hat{X}_{\phi}$, es menor que el nivel de ruido del vacío (establecido como 1/4 en unidades adecuadas). Al medir el espectro de ruido de la fotocorriente diferencia $S(\omega, \phi)$, se puede obtener información sobre estas varianzas. Un valor de $S(\omega, \phi) < 1$ en una banda de frecuencias indica compresión, es decir, ruido cuántico por debajo del límite clásico.

\begin{figure}[h]
    \centering
    \includegraphics[width=0.7\linewidth]{Reporte Semestre 3/imagenes/dhb.png}
    \caption{Esquema de detección homodina balanceada. Un mismo láser impulsa al emisor (círculo rojo) y al oscilador local (OL). Un elemento de fase (no mostrado) permite controlar la fase relativa $\phi$ para seleccionar la cuadratura a medir. La resta de las señales de los dos fotodetectores cancela el ruido clásico del OL.}
    \label{fig:homodina_balanceada}
\end{figure}

\subsection{Ecuación Estocástica de Schrödinger para DHB}

La evolución del estado cuántico condicionado por la medición homodina continua se describe mediante una Ecuación Estocástica de Schrödinger (EES). Para un sistema atómico, la EES en la representación de interacción es \cite{WisemanMilburn2010}:

\begin{eqnarray}
\frac{d|\bar{\Psi}_c\rangle}{dt} = \left[-iH_{eff} + I^W_c(t)\sigma_-\right] |\bar{\Psi}_c\rangle
\end{eqnarray}

donde el Hamiltoniano efectivo $H$ incluye la energía del sistema y la disipación, e $I^W_c(t)$ es la fotocorriente medida, dada por:

\begin{eqnarray}
I^W_c(t) = \left(2\gamma \langle \sigma_{\phi} \rangle_c + \sqrt{\gamma} \, \eta_W(t) \right) e^{-i\phi}
\end{eqnarray}

Aquí, $\gamma$ es la tasa de decaimiento, $\sigma_{\phi} = \frac{1}{2}(\sigma_+ e^{i\phi} + \sigma_- e^{-i\phi})$ es un operador de dipolo dependiente de la fase, $\langle \cdot \rangle_c$ denota el valor esperado condicional, y $\eta_W(t)$ es un ruido blanco gaussiano. Esta EES captura el hecho de que, en la DHB, la medición produce un cambio continuo e infinitesimal en la función de onda, gobernado por un ruido de carácter gaussiano en lugar de los saltos discretos del fotoconteo.


%%% AQUi acaba lo que me voy a traer




\section*{Objetivos}
\textbf{Objetivo General:}

Simular sistemas cuánticos de dos y tres niveles mediante la Teoría de Trayectorias Cuánticas, tanto en su formulación de brincos como en su versión difusiva, con el propósito de estudiar el fenómeno de flouroscencia intermitente, obteniendo como resultados principales la evolución de la función de onda y con ello las poblaciones, la fotocorriente, fotoestadística y espectros.\\

\textbf{Objetivos Específicos:}
\begin{itemize}
    \item Simular un sistema de tres niveles con estructura de shelving bajo detección homodina, analizando cómo la dinámica estocástica se refleja en la fotocorriente y el espectro de fluorescencia en las cuadraturas 0 y pi/2.
    
    \item Desarrollar un análisis fotoestadístico de los tiempos de retorno al estado metaestable aplicando la teoría de procesos de renovación, con el objetivo de caracterizar la distribución de los periodos brillantes y oscuros y validar las predicciones teóricas mediante simulaciones numéricas.
    
    \item Implementar simulaciones numéricas utilizando  ecuaciones estocásticas diferenciales (SDEs) para trayectorias difusivas y trayectorias de saltos.

    \item Generar y analizar una representación de pseudo-esfera de Bloch para el sistema de tres niveles, que permita visualizar la dinámica estocástica del sistema y contrastar cuantitativamente las trayectorias simuladas bajo los enfoques de saltos cuánticos y trayectorias difusivas.
    
    \item Obtener la evolución temporal de las poblaciones y compararlas con las soluciones obtenidas a partir de la ecuación maestra, utilizándolas como referencia para verificar la consistencia y precisión de las simulaciones estocásticas.

\end{itemize}


% Metodología
\section*{Metodología De las Trayectorias de Saltos Cuanticos}

La metodología propuesta para el estudio de los tiempos de espera en sistemas cuánticos de tres niveles se divide en dos partes fundamentales: (1) la simulación de trayectorias estocásticas mediante el método de función de onda condicional de Monte Carlo (MCWF), y (2) la obtención de distribuciones de tiempos de espera para transiciones específicas.

\subsection{Parte 1: Simulación de Trayectorias de Saltos Cuánticos}

\subsubsection{Formulación del Problema}

Consideramos un sistema cuántico de tres niveles con estados $\{|e\rangle, |a\rangle, |g\rangle\}$, donde $|a\rangle$ representa el estado metaestable. La evolución del sistema está gobernada por el Hamiltoniano $H$ y un conjunto de operadores de colapso $\{C_k\}$ que representan los diferentes canales de decaimiento: $C_{eg}$, $C_{ea}$ y $C_{ag}$.

\subsubsection{Algoritmo de Saltos Cuánticos }

El algoritmo implementado sigue la metodología estándar de saltos cuánticos con paso temporal fijo:

\begin{enumerate}
    \item \textbf{Inicialización}: 
    \begin{equation}
        |\Psi_c(0)\rangle = |\psi_{\text{inicial}}\rangle, \quad t = 0
    \end{equation}
    
    \item \textbf{Evolución no-hermítica}:
    La función de onda condicional evoluciona según la ecuación de Schrödinger no-hermítica:
    \begin{equation}
        i\hbar\frac{d}{dt}|\Psi_c\rangle = H_{\text{eff}} |\Psi_c\rangle
        \label{eq:non_hermitian_schrodinger}
    \end{equation}
    donde el Hamiltoniano efectivo es:
    \begin{equation}
        H_{\text{eff}} = H - \frac{i\hbar}{2} \sum_{k=1}^{3} C_k^\dagger C_k
    \end{equation}
    
    \item \textbf{Evolución temporal discreta}:
    Para un paso temporal fijo $\Delta t = 0.01$:
    \begin{equation}
        |\bar{\Psi}_c(t+\Delta t)\rangle = e^{-iH_{\text{eff}}\Delta t/\hbar}|\Psi_c(t)\rangle
        \label{eq:no_emission_evolution}
    \end{equation}
    
    \item \textbf{Cálculo de probabilidades de salto}:
    Las probabilidades de que ocurra un salto en el intervalo $[t, t+\Delta t]$ son:
    \begin{align}
        p_k(t) &= \Delta t \langle\bar{\Psi}_c(t+\Delta t)|C_k^\dagger C_k|\bar{\Psi}_c(t+\Delta t)\rangle \\
        p_{\text{tot}}(t) &= \sum_{k=1}^{3} p_k(t)
    \end{align}
    
    \item \textbf{Decisión estocástica de salto}:
    \begin{itemize}
        \item Se genera un número aleatorio $\xi \sim U(0,1)$
        \item Si $\xi \leq p_{\text{tot}}(t)$: ocurre un salto cuántico
        \item Si $\xi > p_{\text{tot}}(t)$: no hay salto, se actualiza el estado:
        \begin{equation}
            |\Psi_c(t+\Delta t)\rangle = |\bar{\Psi}_c(t+\Delta t)\rangle
        \end{equation}
    \end{itemize}
    
    \item \textbf{Selección del canal de salto}:
    Si ocurre un salto, se selecciona el canal $k$ con probabilidad:
    \begin{equation}
        P(k) = \frac{p_k(t)}{p_{\text{tot}}(t)}
    \end{equation}
    
    \item \textbf{Aplicación del salto cuántico}:
    \begin{equation}
        |\Psi_c(t+\Delta t)\rangle = \frac{C_k |\bar{\Psi}_c(t+\Delta t)\rangle}{\sqrt{\langle\bar{\Psi}_c(t+\Delta t)|C_k^\dagger C_k|\bar{\Psi}_c(t+\Delta t)\rangle}}
    \end{equation}
    
    \item \textbf{Actualización temporal}:
    \begin{equation}
        t \rightarrow t + \Delta t
    \end{equation}
\end{enumerate}

\subsubsection{Reglas de Transición Permitidas}

Las transiciones permitidas dependen del estado actual del sistema:
\begin{itemize}
    \item Desde $|e\rangle$: transiciones permitidas $e \rightarrow g$ y $e \rightarrow a$
    \item Desde $|a\rangle$: transición permitida $a \rightarrow g$ 
    \item Desde $|g\rangle$: no hay transiciones espontáneas
\end{itemize}

\subsubsection{Registro de Trayectorias}

Para cada trayectoria se registran las poblaciones:
\begin{equation}
    \begin{cases}
        \rho_{ee}(t) = |\langle e|\Psi_c(t)\rangle|^2 \\
        \rho_{aa}(t) = |\langle a|\Psi_c(t)\rangle|^2 \\
        \rho_{gg}(t) = |\langle g|\Psi_c(t)\rangle|^2
    \end{cases}
\end{equation}

\subsection*{Parte 2: Obtención de Distribuciones de Tiempos de Espera}

\subsubsection{Definición del Tiempo de Espera}

El tiempo de espera $\tau$ para la transición $e \rightarrow a$ se define como el tiempo hasta el primer salto que corresponde al operador de colapso $C_{ea}$, partiendo del estado inicial $|a\rangle$.

\subsubsection{Protocolo de Simulación Específico}

\begin{enumerate}
    \item \textbf{Condiciones iniciales}: 
    \begin{equation}
        |\Psi_c(0)\rangle = |a\rangle,|g\rangle \text{ o }|e\rangle
    \end{equation}
    
    \item \textbf{Transición objetivo}: 
    \begin{equation}
        \text{Se monitorea específicamente el operador seleccionado } C_{ea},C_{eg}, \text{o } C_{ag}
    \end{equation}
    
    \item \textbf{Criterio de parada}: La simulación se detiene cuando:
    \begin{itemize}
        \item Ocurre la transición seleccionada con éxito
        \item Se alcanza el tiempo máximo $t_{\text{max}}$ (sin transición)
    \end{itemize}
\end{enumerate}

\subsubsection{Implementación Numérica}

\begin{itemize}
    \item \textbf{Paso temporal fijo}: $\Delta t = 0.01$ en todas las simulaciones
    
    \item \textbf{Simulación de múltiples realizaciones}:
    \begin{equation}
        \tau_i = \text{tiempo hasta primer salto} \text{ en la realización } i
    \end{equation}
    
    \item \textbf{Número de realizaciones}: $N \sim 10^3 - 10^5$ para obtener estadísticas confiables
\end{itemize}



\subsection*{Parámetros Numéricos}

Los parámetros clave para la simulación son:
\begin{itemize}
    \item $t_{\text{max}}$: Tiempo máximo de simulación
    \item $\Delta t = 0.01$: Paso temporal fijo
    \item Semillas aleatorias: Para reproducibilidad de resultados
\end{itemize}



Esta metodología permite caracterizar completamente la estadística de los tiempos de espera para transiciones cuánticas específicas, proporcionando información fundamental sobre la dinámica del sistema metaestable.

\section*{Metodología para de las Trayectorias de Detección Homodina}

\subsection*{Evolución difusiva (detección homodina)}

Cuando el entorno se mide mediante detección homodina, el sistema evoluciona de forma continua bajo medición continua de un cuadratura específica del campo. Esta evolución se describe mediante una ecuación de Schrödinger estocástica impulsada por un incremento de Wiener real. 

\subsubsection{Ecuación maestra estocástica}

Para detección homodina, la evolución de la función de onda condicional está dada por:

\begin{equation}
d|\Psi_c\rangle = \left\{-iH_{\text{sys}}dt - \frac{1}{2}\sum_k\left(C_k^\dagger C_k - \langle C_k^\dagger C_k\rangle_c\right)dt + \sum_k\left(C_k - \langle C_k\rangle_c\right)dW_k\right\}|\Psi_c\rangle,
\end{equation}

donde $dW_k$ son incrementos de Wiener reales independientes que satisfacen:

\begin{equation}
\mathbb{E}[dW_k] = 0, \quad dW_k dW_l = \delta_{kl}dt.
\end{equation}

\subsubsection{Hamiltoniano efectivo y operadores de medida}

Para el sistema de tres niveles, el Hamiltoniano efectivo es:

\begin{equation}
H_{\text{eff}} = H_{\text{sys}} - \frac{i}{2}\left(\gamma A_{eg}A_{ge} + \gamma_d A_{ea}A_{ae} + \gamma_a A_{ag}A_{ga}\right),
\end{equation}

con el Hamiltoniano del sistema:

\begin{equation}
H_{\text{sys}} = \Delta A_{ee} + \frac{\Omega}{2}\left(A_{eg} + A_{ge}\right).
\end{equation}

\subsubsection{Fotocorriente homodina}

La fotocorriente medida en detección homodina está dada por:

\begin{equation}
I_c^H(t) = \gamma\langle A_{eg} + A_{ge}\rangle_c + \sqrt{\gamma}\,\eta_H(t),
\end{equation}

donde $\eta_H(t)$ es ruido blanco real que satisface:

\begin{equation}
\mathbb{E}[\eta_H(t)] = 0, \quad \mathbb{E}[\eta_H(t)\eta_H(t')] = \delta(t-t').
\end{equation}

\subsubsection{Algoritmo de simulación para detección homodina}

\begin{enumerate}
\item \textbf{Definir los parámetros del sistema:}
\begin{itemize}
\item $\gamma$: Tasa de emisión espontánea $|e\rangle \to |g\rangle$
\item $\gamma_d$: Tasa de decaimiento $|e\rangle \to |a\rangle$ 
\item $\gamma_a$: Tasa de decaimiento $|a\rangle \to |g\rangle$
\item $\Omega$: Frecuencia de Rabi
\item $\Delta$: Desintonía láser
\end{itemize}

\item \textbf{Definición del espacio de Hilbert y operadores:}

Para el sistema de \textbf{tres niveles} con estados $\{|g\rangle, |e\rangle, |a\rangle\}$, los operadores relevantes son:

\begin{align*}
A_{eg} &= |e\rangle\langle g|, \quad A_{ge} = |g\rangle\langle e|, \\
A_{ea} &= |e\rangle\langle a|, \quad A_{ae} = |a\rangle\langle e|, \\
A_{ag} &= |a\rangle\langle g|, \quad A_{ga} = |g\rangle\langle a|, \\
A_{ee} &= |e\rangle\langle e|, \quad A_{aa} = |a\rangle\langle a|, \quad A_{gg} = |g\rangle\langle g|
\end{align*}

\item \textbf{Operadores de colapso:}
\begin{equation}
C_1 = \sqrt{\gamma}A_{ge}, \quad C_2 = \sqrt{\gamma_d}A_{ae}, \quad C_3 = \sqrt{\gamma_a}A_{ga}
\end{equation}

\item \textbf{Inicialización del estado cuántico:}
\begin{equation}
|\Psi(0)\rangle = |\psi_0\rangle
\end{equation}

\item \textbf{Evolución estocástica:}

Para cada paso de tiempo $dt$:
\begin{enumerate}
    \item Generar incrementos de Wiener $dW_1, dW_2, dW_3$ con distribución normal:
    \begin{equation}
    dW_k \sim \mathcal{N}(0, \sqrt{dt})
    \end{equation}
    
    \item Calcular valores esperados:
    \begin{equation}
    \langle C_k\rangle_c = \langle\Psi_c|C_k|\Psi_c\rangle
    \end{equation}
    
    \item Evolucionar la función de onda:
    \begin{align*}
    d|\Psi_c\rangle = & \left[-iH_{\text{sys}}dt - \frac{1}{2}\sum_{k=1}^3\left(C_k^\dagger C_k - \langle C_k^\dagger C_k\rangle_c\right)dt\right. \\
    & \left. + \sum_{k=1}^3\left(C_k - \langle C_k\rangle_c\right)dW_k\right]|\Psi_c\rangle
    \end{align*}
    
    \item Normalizar el estado:
    \begin{equation}
    |\Psi_c(t+dt)\rangle = \frac{|\Psi_c(t) + d|\Psi_c\rangle}{\||\Psi_c(t) + d|\Psi_c\rangle\|}
    \end{equation}
    
    \item Calcular la fotocorriente homodina:
    \begin{equation}
    I_c^H(t) = \gamma\langle A_{eg} + A_{ge}\rangle_c + \sqrt{\gamma}\,\frac{dW_1}{\sqrt{dt}}
    \end{equation}
\end{enumerate}

\item \textbf{Cálculo de correlaciones y espectro:}

La función de correlación de la fotocorriente está dada por:
\begin{equation}
G^{(2)}(\tau) = \mathbb{E}[I_c^H(t+\tau)I_c^H(t)]
\end{equation}

El espectro de emisión se obtiene mediante:
\begin{equation}
S(\omega) = 1 + \gamma\int_0^\infty d\tau\,\cos(\omega\tau)\langle A_{eg}(t+\tau)A_{ge}(t)\rangle
\end{equation}

\item \textbf{Promedio sobre realizaciones:}

Para obtener resultados estadísticamente significativos, se promedian múltiples realizaciones del proceso estocástico:
\begin{equation}
\langle O\rangle = \frac{1}{N}\sum_{i=1}^N \langle\Psi_c^i|O|\Psi_c^i\rangle
\end{equation}
\end{enumerate}

\subsubsection{Consideraciones numéricas}

\begin{itemize}
\item \textbf{Paso temporal}: $dt = 0.01$ para balancear precisión y eficiencia computacional
\item \textbf{Número de realizaciones}: $N \sim 10^3 - 10^4$ para convergencia estadística
\item \textbf{Tiempo de simulación}: Suficientemente largo para capturar la dinámica característica
\item \textbf{Esquema de integración}: Euler-Maruyama para la evolución estocástica
\end{itemize}

\subsubsection{Implementación práctica}

La implementación numérica sigue estos pasos:

\begin{enumerate}
\item Inicializar el estado $|\Psi_c(0)\rangle$ y arrays de almacenamiento
\item Para cada paso de tiempo:
\begin{enumerate}
    \item Calcular todos los valores esperados necesarios
    \item Generar ruido estocástico
    \item Aplicar la evolución difusiva
    \item Normalizar el estado
    \item Almacenar observables de interés
\end{enumerate}
\item Repetir para múltiples realizaciones con diferentes semillas aleatorias
\item Promediar sobre realizaciones para obtener valores esperados
\item Calcular funciones de correlación y espectros
\end{enumerate}

Esta metodología permite estudiar la dinámica cuántica bajo medición continua homodina y caracterizar completamente las propiedades estadísticas del sistema de tres niveles.
% Resultados
% En el documento:
\section*{Resultados}

\subsection*{1.- Distribuciones de Tiempos Oscuros y Brillantes}

La caracterización completa del la Florescencia Intermitente requiere verificar las distribuciones individuales de los tiempos oscuros y brillantes. La figura \ref{fig:dark_bright_times} muestra estas distribuciones para $\Omega = 3.5\gamma$, donde se observa el comportamiento exponencial característico predicho por la teoría.

\begin{figure}[H]
\centering
\includegraphics[width=1\textwidth]{Reporte Semestre 3/imagenes/Duracion peridos bo.png}
\caption{Distribuciones de (a) tiempos oscuros y (b) tiempos brillantes para $\Omega = 3.5\gamma$. Las líneas negras representan las curvas teóricas $P_O(t) = \frac{1}{T_O}e^{-t/T_O}$ y $P_B(t) = \frac{1}{T_B}e^{-t/T_B}$, mostrando que concuerdan con los datos de simulación.}
\label{fig:dark_bright_times}
\end{figure}

Los resultados numéricos confirman las predicciones analíticas:

\begin{itemize}
\item \textbf{Tiempos oscuros:} $T_O^{teo} = 66.7$ vs $T_O^{med} = 67.8$ (error del 1.7\%)
\item \textbf{Tiempos brillantes:} $T_B^{teo} = 41.5$ vs $T_B^{med} = 41.9$ (error del 1\%)
\end{itemize}

La naturaleza exponencial de ambas distribuciones valida el modelo de procesos de Poisson independientes para las transiciones hacia y desde el estado metaestable. La calidad del ajuste demuestra que la simulación captura correctamente la física fundamental del proceso de intermitencia.



\subsection*{2.- Distribuciones de Tiempos de Retorno al Estado Metaestable}

Para caracterizar la dinámica de retorno al estado metaestable $|a\rangle$, implementamos el algoritmo de saltos cuánticos con las siguientes condiciones específicas:

\begin{itemize}
    \item \textbf{Estado inicial}: El sistema se prepara en el estado metaestable 
    \[
    |\Psi(0)\rangle = |a\rangle
    \]
    
    \item \textbf{Transición monitorizada}: Se registra el tiempo hasta la primera ocurrencia de la transición $e \rightarrow a$, correspondiente al operador de colapso $C_{ea}$
    
    \item \textbf{Criterio de parada}: Cada realización de la simulación finaliza cuando:
    \begin{itemize}
        \item Ocurre la transición $C_{ea}$ (éxito)
        \item Se alcanza el tiempo máximo de simulación $t_{\text{max}}$ (sin transición)
    \end{itemize}
\end{itemize}

Esta configuración permite obtener una muestra estadística de los tiempos de retorno $\{\tau_i\}$ al estado metaestable, a partir de la cual se construye la distribución de probabilidad $P(\tau)$ que caracteriza la dinámica de repoblación del nivel $|a\rangle$.
\begin{figure}[H]
    \centering
    \includegraphics[width=0.5\linewidth]{Reporte Semestre 1/1 sistema.png}
    \caption{Esquema del sistema de 3 niveles}
    \label{fig:placeholder}
\end{figure}
Para validar nuestro formalismo analítico, realizamos simulaciones numéricas de las trayectorias cuánticas y comparamos los resultados con las predicciones teóricas. La figura \ref{fig:return_times} muestra la distribución de tiempos de retorno al estado metaestable para dos frecuencias de Rabi distintas, donde se observa un excelente acuerdo entre la teoría (línea continua) y la simulación (histograma).

\begin{figure}[H]
\centering
\includegraphics[width=1\textwidth]{Reporte Semestre 3/retorno.png}
\caption{Distribución de tiempos de retorno al estado metaestable para (izquierda) $\Omega = 0.54\gamma$ y (derecha) $\Omega = 3.5\gamma$. La curva roja/azul representa la distribución teórica $P_{OB}(t) = \frac{1}{T_B - T_O}\left(e^{-t/T_B} - e^{-t/T_O}\right)$, mientras que el histograma muestra los datos de la simulación.}
\label{fig:return_times}
\end{figure}

Los resultados numéricos confirman nuestras predicciones analíticas con una precisión notable:

\begin{itemize}
\item \textbf{Para $\Omega = 3.5\gamma$:} El tiempo promedio medido ($T_{OB}^{med} = 108.995$) coincide casi perfectamente con el valor teórico ($T_{OB}^{teo} = 108.140$), presentando un error de solo \textbf{0.8 porciento}. Esta muy buena aproximación valida nuestro formalismo en el régimen de Rabi fuerte.

\item \textbf{Para $\Omega = 0.54\gamma$:} Aunque el acuerdo sigue siendo muy bueno ($T_{OB}^{med} = 180.607$ vs $T_{OB}^{teo} = 168.567$, error del \textbf{7.1\%}), la pequeña discrepancia se atribuye a efectos estadísticos en las colas de la distribución, donde se requieren trayectorias más largas para una muestreo adecuado.

\end{itemize}

\textbf{Análisis físico del comportamiento observado:}

La marcada diferencia en las distribuciones para distintos $\Omega$ revela física importante:

\begin{itemize}
\item \textbf{Efecto de la frecuencia de Rabi:} En el caso $\Omega = 0.54\gamma$ (baja frecuencia de Rabi), el periodo brillante se alarga significativamente ($T_B = 101.9$ vs $T_B = 41.5$ para $\Omega = 3.5\gamma$). Esto ocurre porque a menor $\Omega$, el sistema tarda más en excitarse desde el estado base al excitado, prolongando así los periodos de emisión.

\item \textbf{Colas de la distribución:} La distribución para $\Omega = 0.54\gamma$ exhibe una cola más extendida, indicando la presencia de ciclos excepcionalmente largos. Estos eventos raros contribuyen al mayor tiempo promedio observado en este régimen.

\item \textbf{Simetría de la distribución:} Para $\Omega = 3.5\gamma$, la distribución es más simétrica y concentrada alrededor del promedio, reflejando la mayor eficiencia del bombeo óptico y ciclos de parpadeo más regulares.

\end{itemize}

La consistencia entre teoría y simulación, particularmente el error inferior al 1 porciento para $\Omega = 3.5\gamma$, demuestra la robustez de nuestro enfoque analítico y valida la interpretación física del Florescencia Resonante Intermitente como un proceso de dos escalas temporales acopladas.




\subsection*{Espectroscopía Heterodina de un Emisor Cúantico}

Mejore las gráficas de deteccion heterodina aumentando la resolución temporal y el tiempo de simulación obteniendo un mejor espectro que muestra el pico central angosto 

\subsubsection*{Validación de las Trayectorias Cuánticas}

La figura \ref{fig:heterodyne_convergence} muestra la convergencia entre el promedio sobre 500 trayectorias heterodinas y la solución de la ecuación maestra para la población del estado excitado.Esto valida nuestra implementación de las ecuaciones estocásticas y demuestra la robustez del método de detección heterodina para reproducir correctamente la dinámica coherente del sistema.

\begin{figure}[h!]
\centering
\includegraphics[width=0.8\textwidth]{Reporte Semestre 3/imagenes/convergencia.png}
\caption{Convergencia entre el promedio sobre 500 trayectorias heterodinas (línea azul) y la solución de la ecuación maestra (línea naranja punteada) para la población del estado excitado. La coincidencia valida la implementación numérica del esquema de detección heterodina.}
\label{fig:heterodyne_convergence}
\end{figure}

\subsubsection*{Estructura Espectral y el Triplete de Mollow}

La figura \ref{fig:mollow_triplet} presenta los espectros de emisión para dos regímenes distintos de frecuencia de Rabi, obtenidos mediante detección heterodina con un tiempo de simulación de $2000\gamma t$ y una resolución temporal de $0.01\gamma t$.

\begin{figure}[h!]
\centering
\includegraphics[width=\textwidth]{Reporte Semestre 3/imagenes/dhete.png}
\caption{Espectros de emisión por detección heterodina. (Izquierda) Para $\Omega = 3.5\gamma$ se observa el triplete de Mollow característico, con el pico central angosto que manifiesta el fenómeno de fluorescencia intermitente a través del decaimiento lento de las coherencias. (Derecha) Para $\Omega = 0.54\gamma$ (frecuencia de Rabi óptima) se obtiene el pico central más angosto posible con desintonía $\Delta = 0$, destacando el régimen de emisión resonante donde la fluorescencia intermitente es más pronunciada.}
\label{fig:mollow_triplet}
\end{figure}

En el espectro para $\Omega = 3.5\gamma$ (panel izquierdo) se identifica claramente el triplete de Mollow, consistente con:
\begin{itemize}
\item \textbf{Dos picos laterales} en $\omega_0 \pm \Omega$, correspondientes a las transiciones laterales del dressing atómico inducido por el campo de bombeo.
\item \textbf{Un pico central angosto} en $\omega_0$, cuya anchura reducida es consecuencia directa del fenómeno de fluorescencia intermitente. Este estrechamiento espectral refleja los largos tiempos de atrapamiento en el estado metaestable $|a\rangle$, que generan períodos extendidos sin emisión de fotones y, por tanto, una componente espectral de baja frecuencia muy definida.
\end{itemize}

El panel derecho ($\Omega = 0.54\gamma$) revela un comportamiento fundamental: esta frecuencia de Rabi específica maximiza el efecto de fluorescencia intermitente, produciendo el pico central más angosto achievable en resonancia ($\Delta = 0$). La optimización de $\Omega$ permite sintonizar la dinámica de atrapamiento en el estado metaestable, lo que se manifiesta espectroscópicamente como un estrechamiento extremo del pico central debido a la prolongación de los intervalos sin emisión.

\subsection*{Detección Homodina y Cuadraturas del Campo}

Para complementar el análisis espectral y acceder a las propiedades de fase del campo emitido, implementamos un esquema de detección homodina. En esta configuración, la señal del emisor se mezcla con un oscilador local de la misma frecuencia, permitiendo medir las cuadraturas específicas del campo electromagnético mediante el ajuste del ángulo de fase $\phi$ del oscilador local.

\subsubsection*{Análisis por Cuadraturas}

La figura \ref{fig:homodyne_quadratures} muestra los espectros homodinos para dos cuadraturas ortogonales, obtenidos con $\phi = 0$ y $\phi = \pi/2$, simulados bajo las mismas condiciones de tiempo ($2000\gamma t$) y resolución temporal ($\Delta t=0.01$) que el caso heterodino.

\begin{figure}[h!]
\centering
\includegraphics[width=\textwidth]{Reporte Semestre 3/imagenes/dhomodinaaa.png}
\caption{Espectros de detección homodina para diferentes cuadraturas. (Izquierda) Cuadratura $\phi = 0$: revela exclusivamente la componente coherente del pico central, suprimiendo los picos laterales. (Derecha) Cuadratura $\phi = \pi/2$: muestra los picos laterales del triplete de Mollow y el pico central angosto asociado al decaimiento lento hacia el estado metaestable.}
\label{fig:homodyne_quadratures}
\end{figure}

Los resultados demuestran la selectividad de fase inherente a la detección homodina:

\begin{itemize}
\item \textbf{Cuadratura $\phi = 0$:} Resulta particularmente reveladora, ya que exhibe \textbf{únicamente la componente del pico central}, suprimiendo completamente los picos laterales del triplete de Mollow.

\item \textbf{Cuadratura $\phi = \pi/2$:} Esta configuración revela la parte faltante del triplete de Mollow, los picos laterales ademas muestra un pico central angosto. La anchura reducida de este pico central en $\phi = \pi/2$ refleja directamente la dinámica lenta de decaimiento hacia el estado metaestable $|a\rangle$, que introduce una componente espectral de baja frecuencia muy definida.
\end{itemize}

\subsubsection*{Implicaciones Físicas}

La capacidad de aislar selectivamente diferentes componentes espectrales mediante la elección de la cuadratura tiene importantes consecuencias:

\begin{itemize}
\item La supresión completa de los picos laterales en $\phi = 0$ demuestra que estas características espectrales corresponden exclusivamente a emisión en la cuadratura $\phi = \pi/2$.
\item El pico central angosto presente en $\phi = \pi/2$ manifiesta la dinámica de decaimiento lento hacia el estado metaestable, donde los largos tiempos característicos de esta transición se traducen en un estrechamiento espectral significativo.
\item Esta separación de cuadraturas revela que los picos laterales del triplete de Mollow son puramente incoherentes y aparecen solo en una cuadratura específica, mientras que la emisión coherente se manifiesta predominantemente en la cuadratura ortogonal.
\item La presencia del pico central angosto en $\phi = \pi/2$ proporciona evidencia espectroscópica directa de la existencia del estado metaestable y su influencia en la dinámica temporal del sistema.
\end{itemize}


\subsection*{Dinámica Cuántica en la Pseudo-Esfera de Bloch}

Para visualizar directamente la evolución temporal del sistema cuántico y contrastar los diferentes regímenes de medición, construimos una representación en pseudo-esfera de Bloch donde los ejes corresponden a las poblaciones normalizadas de los estados $|e\rangle$, $|g\rangle$ y $|a\rangle$ ($P_{ee}$, $P_{gg}$, $P_{aa}$). Esta representación tridimensional nos permite seguir la trayectoria del estado cuántico en tiempo real bajo dos esquemas de medición fundamentalmente diferentes.

\subsubsection*{Dinámica con Saltos Cuánticos}

La figura \ref{fig:jump_dynamics} muestra la evolución del sistema bajo un esquema de medición que registra saltos cuánticos discretos. Las características principales observadas son:

\begin{itemize}
\item \textbf{Movimiento suave entre saltos:} La trayectoria evoluciona de manera continua y determinista entre los eventos de salto, reflejando la evolución unitaria del sistema bajo el Hamiltoniano efectivo.


\item \textbf{Transiciones abruptas:} Los saltos cuánticos aparecen como cambios discontinuos en la posición sobre la esfera, correspondientes a la detección de fotones o transiciones al estado metaestable.

\item \textbf{Visitas al estado metaestable:} Ocasionalmente se observan excursiones hacia la región correspondiente al estado $|a\rangle$, donde el sistema permanece atrapado temporalmente, manifestando el fenómeno de intermitencia.
\end{itemize}

\begin{figure}[H]
\centering
\includegraphics[width=0.8\textwidth]{Reporte Semestre 3/imagenes/pseodoesfera tjsaltos.png}
\caption{Trayectoria en la pseudo-esfera de Bloch para dinámica con saltos cuánticos. Los segmentos suaves corresponden a evolución entre saltos, mientras que los cambios abruptos indican detección de fotones o transiciones al estado metaestable.}
\label{fig:jump_dynamics}
\end{figure}

\subsubsection*{Dinámica Difusiva}

La figura \ref{fig:diffusive_dynamics} presenta el comportamiento bajo un esquema de medición homodina que genera una evolución difusiva continua:

\begin{itemize}
\item \textbf{Trayectoria continua:} A diferencia del caso con saltos, aquí la evolución es completamente continua, sin discontinuidades abruptas.

text
\item \textbf{Difusión suave:} El estado cuántico exhibe un movimiento browniano sobre la esfera, reflejando el ruido de medición continuo inherente a la detección homodina.

\item \textbf{Ausencia de visitas abruptas al metaestable:} La transición hacia el estado $|a\rangle$ ocurre de manera gradual, sin los saltos discretos característicos del régimen anterior.
\end{itemize}

\begin{figure}[H]
\centering
\includegraphics[width=0.8\textwidth]{Reporte Semestre 3/imagenes/Esfera difusiva.png}
\caption{Trayectoria difusiva en la pseudo-esfera de Bloch para medición homodina. La evolución es completamente continua, mostrando un comportamiento browniano que refleja el ruido de medición.}
\label{fig:diffusive_dynamics}
\end{figure}


\newpage
%Trabajo Futuro
\section*{Perspectiva del Cuarto Semestre}
Durante el cuarto semestre, se abordarán diferentes líneas de investigación y formación académica con el objetivo de profundizar en fenómenos cuánticos avanzados y técnicas de detección en sistemas de múltiples niveles. A continuación, se detallan los principales puntos de enfoque:
\begin{enumerate}

\item Detección Homodina Condicionada y \textit{squeezing}: 
Se explorará un marco teórico aún no desarrollado dentro del proyecto, enfocado en cómo mediciones homodinas condicionadas pueden modificar dinámicas cuánticas y dar lugar a estados comprimidos de la luz. Este estudio permitirá evaluar si dichas técnicas pueden incorporarse a los modelos numéricos existentes y ofrecer nuevas rutas para la caracterización de fluctuaciones cuánticas.

\item Tópicos de Electrónica Fotónica y Cómputo: Espectroscopía Láser: 
Como parte de la última asignatura de la maestría, se profundizará en técnicas modernas de espectroscopía láser relevantes para el estudio de sistemas atómicos abiertos. 

\item Colaboración en simulación de Batimientos Cu\'anticos en un sistema de cuatro niveles con Detección Heterodina:
Se propone utilizar la detección heterodina para simular el sistema de cuatro niveles donde aparecen fenómenos de quantum beats. Esta simulación se basará en resultados previos obtenidos en el grupo de investigación.%, los cuales proporcionan un punto de partida sólido para la implementación numérica. 

% El curso proporcionará herramientas teóricas y prácticas necesarias para el desarrollo de experimentos y simulaciones que involucren fenómenos ópticos cuánticos de alta precisión.

\item Colaboración en Estudios Radiaci\'on Cooperativa. Existe la posibilidad de aportar en este proyecto como una línea de trabajo paralela a la tesis principal. En particular, se colaborará en el estudio teórico y numérico de sistemas compuestos por dos átomos interactuando, donde pueden emerger fenómenos de \textit{subradiancia} y \textit{superradiancia}.  En dichos sistemas, es posible observar picos angostos en el espectro de emisión, los cuales están asociados a la supresión o el refuerzo colectivo de la radiación. 





\end{enumerate}
\newpage
\section*{Anexos}


\subsection*{Cronograma de actividades}


\begin{table}[h!]
\centering
\renewcommand{\arraystretch}{1.5}
\begin{tabular}{|p{9cm}|c|c|c|c|c|}
\hline
\textbf{Actividad} & \multicolumn{5}{c|}{\textbf{Semestre}} \\ \cline{2-6} 
 & 1 & 2 & 3 & 4 & 5 \\ \hline

Búsqueda de información y análisis del estado del Arte & \cellcolor{green} & & & & \\ \hline
Estudio detallado de ecuaciones maestras y teoría de trayectorias cuánticas & \cellcolor{green} & & & & \\ \hline
Definición del marco teórico y planteamiento del problema & \cellcolor{green} & & & & \\ \hline

Definición de la metodología de simulación numérica (trayectorias cuánticas) & & \cellcolor{green} & & & \\ \hline
Selección y adaptación de herramientas computacionales (QuTiP, Python) & & \cellcolor{green} & & & \\ \hline
Elaboración del diseño de simulación (condiciones iniciales, parámetros, algoritmos) & & \cellcolor{green} & & & \\ \hline
Programación de simulaciones de trayectorias cuánticas para sistemas simples & & \cellcolor{green} & & & \\ \hline
Simulación y ajuste de parámetros (ejecución y depuración de código) & & \cellcolor{green} & & & \\ \hline
Análisis preliminar de resultados de simulaciones & & \cellcolor{green} & & & \\ \hline

Comparación con experimentos previos y ecuaciones maestras & & & \cellcolor{green} & & \\ \hline
Validación y ajuste de modelos teóricos & & & \cellcolor{green} & & \\ \hline
Discusión de resultados, refinamiento de hipótesis & & & \cellcolor{green} & & \\ \hline

Redacción del informe de investigación (marco teórico, metodología, resultados) & & & \cellcolor{green} & \cellcolor{green} & \cellcolor{green} \\ \hline
Preparación de presentaciones para congresos o seminarios & & & & \cellcolor{green} & \cellcolor{green} \\ \hline
Defensa del proyecto & & & & \cellcolor{green} & \cellcolor{green} \\ \hline

\end{tabular}
\end{table}



\newpage
\begin{thebibliography}{99}

\bibitem{Carmichael1993} 
Carmichael, H. J. (1993). \textit{An Open Systems Approach to Quantum Optics}. Springer.

\bibitem{Plenio1998} 
Plenio, M. B., \& Knight, P. L. (1998). The quantum jump approach to dissipative dynamics. \textit{Rev. Mod. Phys.}, 70(1), 101-144.

\bibitem{Knight1990} 
Knight, P. L. (1990). Quantum fluctuations in dissipative systems. In \textit{Quantum Noise in Semiconductor Lasers} (pp. 1-23). Elsevier.

\bibitem{Wiseman1993} 
Wiseman, H. M., \& Milburn, G. J. (1993). Interpretation of quantum jump and diffusion processes illustrated on the Bloch sphere. \textit{Phys. Rev. A}, 47(3), 1652-1666.

\bibitem{WisemanMilburn2010}
Wiseman, H. M., \& Milburn, G. J. (2010). \textit{Quantum Measurement and Control}. Cambridge University Press.

\bibitem{Wiseman1996}
Wiseman, H. M. (1996). Quantum trajectories and quantum measurement theory. \textit{Quantum and Semiclassical Optics}, 8(1), 205-222.

\bibitem{Daley2014}
Daley, A. J. (2014). Quantum trajectories and open many-body systems. \textit{Advances in Physics}, 63(2), 77-149.

\bibitem{Lindner2022}
Lindner, N., \& Banerjee, D. (2022). Efficient tensor-network approach for quantum trajectories. \textit{Physical Review Letters}, 128, 150605.

\bibitem{Li2022}
Li, J., et al. (2022). GPU-accelerated quantum dynamics with QuTiP. \textit{Quantum}, 6, 630.

\bibitem{Koch2019}
Koch, J., \& Girvin, S. M. (2019). Introduction to circuit QED: Photon-assisted dephasing and quantum trajectories. \textit{Annalen der Physik}, 525(6), 395-416.

\bibitem{Suess2014}
Suess, D., Eisfeld, A., \& Strunz, W. T. (2014). Hierarchy of stochastic pure states for open quantum system dynamics. \textit{Physical Review Letters}, 113, 150403.

\bibitem{Ke2021}
Ke, Y., Tang, J., Xu, R., \& Yan, Y. (2021). A stochastic Schrödinger equation approach to quantum dissipative dynamics. \textit{The Journal of Chemical Physics}, 154(8), 084103.

\bibitem{Minev2019}
Minev, Z. K., et al. (2019). To catch and reverse a quantum jump. \textit{Nature}, 570, 200-204.

\bibitem{Haug2023}
Haug, T., Kim, M. S., \& Paternostro, M. (2023). Variational quantum trajectory method for open quantum systems. \textit{npj Quantum Information}, 9, 20.

\bibitem{Vicentini2022}
Vicentini, F., et al. (2022). Neural-network quantum states for open quantum systems. \textit{Physical Review Letters}, 129, 220602.

\end{thebibliography}


\end{document}
\end{document}